% Figure: converging-diverging (de Laval) nozzle schematic with the
% Mach-number progression along the axis. Subsonic in the converging
% section, M=1 at the throat, supersonic in the diverging section.
\begin{figure}[htbp]
\centering
\begin{tikzpicture}[scale=1.0]
% nozzle wall, upper and lower (mirror); converging then diverging
% profile y(x): inlet half-height 1.6, throat 0.5, exit 1.3
\def\inlet{1.6}
\def\throat{0.5}
\def\exit{1.3}
% upper wall as a smooth-ish polyline
\draw[engblue, very thick]
  (0,\inlet) -- (1.0,1.3) -- (2.0,0.9) -- (3.0,\throat)
  -- (4.0,0.85) -- (5.0,1.1) -- (6.0,\exit);
% lower wall (mirror)
\draw[engblue, very thick]
  (0,-\inlet) -- (1.0,-1.3) -- (2.0,-0.9) -- (3.0,-\throat)
  -- (4.0,-0.85) -- (5.0,-1.1) -- (6.0,-\exit);
% throat dashed line
\draw[enggray, dashed] (3.0,-\throat) -- (3.0,\throat);
\node[enggray, font=\scriptsize, anchor=south] at (3.0,0.55) {throat $A^{*}$, $M=1$};
% centreline with flow arrow
\draw[engblack, -{Latex[length=2.5mm]}] (-0.4,0) -- (6.4,0);
% region labels
\node[engblack, font=\scriptsize] at (1.3,-0.15) {$M<1$};
\node[engred, font=\scriptsize] at (4.6,-0.15) {$M>1$};
\node[engblack, font=\scriptsize, anchor=east] at (-0.45,0.9) {$p_0,T_0$};
% inset Mach profile below
\begin{scope}[yshift=-3.6cm]
  \draw[engblack, -{Latex[length=2mm]}] (-0.4,0) -- (6.4,0) node[right, font=\scriptsize] {$x$};
  \draw[engblack, -{Latex[length=2mm]}] (-0.4,0) -- (-0.4,2.2) node[above, font=\scriptsize] {$M$};
  % M=1 reference
  \draw[enggray, dashed] (-0.4,1.0) -- (6.4,1.0);
  \node[enggray, font=\scriptsize, anchor=east] at (-0.45,1.0) {$1$};
  % Mach curve: rising from ~0.2 to 1 at throat (x=3) then to ~2.4 at exit
  \draw[engred, very thick] plot[smooth] coordinates {
    (0,0.18) (1.0,0.32) (2.0,0.6) (3.0,1.0)
    (4.0,1.55) (5.0,2.0) (6.0,2.35)};
  \draw[enggray, dashed] (3.0,0) -- (3.0,1.0);
\end{scope}
\end{tikzpicture}
\caption[Converging-diverging nozzle and the Mach progression]{The
converging-diverging (de Laval) nozzle. The subsonic inlet flow
accelerates through the converging section, reaches $M=1$ at the
throat (the minimum area $A^{*}$), and accelerates supersonically
through the diverging section. The lower panel plots the centreline
Mach number against axial position for the design (fully supersonic)
operating mode: it rises monotonically through the throat and beyond.
The divergence after the throat is required because the density falls
faster than the velocity rises once the flow is supersonic, so the
area must grow to conserve mass. This is the geometry of every rocket
nozzle and supersonic wind-tunnel nozzle.}
\label{fig:vol03ch13:cd-nozzle}
\end{figure}
